\phantomsection
\section*{修订历史}\label{revision-history}
\lhead{\it{修订历史}}
\rhead{\Navigation}
\addcontentsline{toc}{chapter}{修订历史}

\begin{longtable}[l]{ | m{2cm} | m{1.5cm} | m{12cm} | }\hline

    \rowcolor{lightgray}
    日期 & 版本 & \hypertarget{Release notes}{发布说明}\\\hline
    \endfirsthead

    \multicolumn{3}{c}{\bfseries 续上页}\\\hline
    \rowcolor{lightgray}
    日期 & 版本 & 发布说明\\\hline
    \endhead

    \multicolumn{3}{r}{\textbf{见下页}}\\
    \endfoot

    \endlastfoot

    % table contents

2026 & v5.8 &
\vspace{0.5em}
\begin{itemize}
    \item 章节 \ref{mod:sensor} \textit{\nameref{mod:sensor}}：新增寄存器 \hyperref[regdesc:SENSSARMEASCTRL2REG]{SENS\_SAR\_MEAS\_CTRL2\_REG} %DOC-6803
\end{itemize}
\vspace{-1em}
\\\hline
2026.03 & v5.7 &
\vspace{0.5em}
\begin{itemize}
    \item 章节 \ref{mod:ethernetmac} \textit{\nameref{mod:ethernetmac}}：更新寄存器 \hyperref[regdesc:DMAMISSEDFRREG]{DMAMISSEDFR\_REG} 的 bit[31:29] 为保留位，并更新域 Missed\_FC 为 bit[15:0] %DOC-11142
\end{itemize}
\vspace{-1em}
\\\hline
2025.11 & v5.6 &
\vspace{0.5em}
\begin{itemize}
    \item 章节 \ref{mod:lowpowm} \textit{\nameref{mod:lowpowm}}：更正拒绝睡眠源的描述
\end{itemize}
\vspace{-1em}
\\\hline

2025.08 & v5.5 &
\vspace{0.5em}
\begin{itemize}
    \item 章节 \ref{mod:i2c} \textit{\nameref{mod:i2c}}：
        \begin{itemize}
            \item 在寄存器 \hyperref [regdesc:I2CFIFOCONFREG]{I2C\_FIFO\_CONF\_REG} 中增加 bit[9:0]，在寄存器 \hyperref [regdesc:I2CINTRAWREG]{I2C\_INT\_RAW\_REG}、\hyperref [regdesc:I2CINTCLRREG]{I2C\_INT\_CLR\_REG}、\hyperref [regdesc:I2CINTENAREG]{I2C\_INT\_ENA\_REG}、\hyperref [regdesc:I2CINTSTATUSREG]{I2C\_INT\_STATUS\_REG} 中增加 bit[2:0] 和 bit4
            \item 在章节 \ref{sec:i2c-reg-summ} \textit{\nameref{sec:i2c-reg-summ}} 中增加寄存器 \hyperref[regdesc:I2CCTRREG]{I2C\_CTR\_REG}、\hyperref[regdesc:I2CSRREG]{I2C\_SR\_REG} 和 \hyperref[regdesc:I2CTOREG]{I2C\_TO\_REG}
        \end{itemize}
    \item 章节 \ref{mod:sdioslave} \textit{\nameref{mod:sdioslave}}：在寄存器 \hyperref[regdesc:SLCCONF1SLCCONF1REG]{SLCCONF1\_REG} 中增加 bit[3:0]
    \item 章节 \ref{mod:ulp} \textit{\nameref{mod:ulp}}：更新 SLEEP 指令中 sleep\_reg 的比特位
\end{itemize}
\vspace{-1em}
\\\hline


2025.05 & v5.4 &
\vspace{0.5em}
 \begin{itemize}
        \item 章节 \ref{mod:dportreg} \textit{\nameref{mod:dportreg}}：在 \ref{sec:dport-reg-summ} \textit{\nameref{sec:dport-reg-summ}} 小节中增加访问异常寄存器及描述
        \item 章节 \ref{mod:ethernetmac} \textit{\nameref{mod:ethernetmac}}：修正寄存器描述小节中一处笔误
        \item 章节 \ref{mod:uart} \textit{\nameref{mod:uart}}：更新图 \ref{Figure:uart_mem} \textit{\nameref{Figure:uart_mem}}
\end{itemize}
\vspace{-1em}
\\\hline

2025.01 & v5.3 &
\vspace{0.5em}
 \begin{itemize}
        \item 章节 \ref{mod:iomuxgpio} \textit{\nameref{mod:iomuxgpio}}：
        \begin{itemize}
            \item 在 \ref{fig:pad-overview} \textit{\nameref{fig:pad-overview}} 小节中新增一条脚注 %DOC-9726
            \item 新增图 \ref{fig:pad-overview} \textit{\nameref{fig:pad-overview}} %DOC-9726
            \item 更新 \ref{sec:peripheral_signals} \textit{\nameref{sec:peripheral_signals}} 小节 %doc-8358
            \item 更新 \ref{section:4.12} \textit{\nameref{section:4.12}} 小节和 \ref{section:4.13} \textit{\nameref{section:4.13}} 小节，仅显示用户可以用的 GPIO %DOC-9413 %DOC-7678
        \end{itemize}
    \item 章节 \ref{mod:i2c} \textit{\nameref{mod:i2c}}：在小节 \ref{I2C Architecture} 里更新有关 NON-FIFO 模式的描述
  \end{itemize}
  \vspace{-1em}
  \\\hline

2024.08 & v5.2 &
\vspace{0.5em}
 \begin{itemize}
    \item 章节 \ref{mod:ethernetmac} \textit{\nameref{mod:ethernetmac}}：在小节 \ref{emac-rmii-interface-signal} 里增加关于参考时钟的限制 %doc-8456
     \item 章节 \ref{mod:pcnt} \textit{\nameref{mod:pcnt}}：更正 PCNT\_U\textit{\textcolor{red}{\it{n}}}\_CNT\_REG (\textcolor{red}{\it{n}}: 0-7) 的地址 % DOC-7654
     \item 章节 \ref{mod:ledpwm} \textit{\nameref{mod:ledpwm}}：更正部分寄存器的地址 % DOC-7653
     \item 章节 \ref{mod:lowpowm} \textit{\nameref{mod:lowpowm}}：
        \begin{itemize}
            \item 更新有关预设功耗模式的描述 %TRMC6-42193
            \item 增加表 \ref{tab:lpm-wakeup-source} 下方有关 EXT1 的说明 %DOC-8414
        \end{itemize}
     \item 章节 \ref{mod:iomuxgpio} \textit{\nameref{mod:iomuxgpio}}：
     \begin{itemize}
         \item 更新以下寄存器 \hyperref[regdesc:GPIOSTATUSREG]{GPIO\_STATUS\_REG}、\hyperref[regdesc:GPIOSTATUSW1TSREG]{GPIO\_STATUS\_W1TS\_REG}、\hyperref[regdesc:GPIOSTATUSW1TCREG]{GPIO\_STATUS\_W1TC\_REG}、\hyperref[regdesc:GPIOSTATUS1REG]{GPIO\_STATUS1\_REG}、\hyperref[regdesc:GPIOSTATUS1W1TSREG]{GPIO\_STATUS1\_W1TS\_REG} 和 \hyperref[regdesc:GPIOSTATUS1W1TCREG]{GPIO\_STATUS1\_W1TC\_REG} 的描述 %DOC-8148
         \item 更新章节 \ref{PadHold} 中对 Hold 控制的描述 % DOC-7326
     \end{itemize}
 \end{itemize}
 \vspace{-1em}
\\\hline

2024.04 & v5.1 &
\vspace{0.5em}
\begin{itemize}

    \item 章节 \ref{mod:iomuxgpio} \textit{\nameref{mod:iomuxgpio}}：
    \begin{itemize}
        \item 更新 \ref{PadHold} 中的描述 %DOC-5892
        \item 更新以下寄存器的描述：RTCIO\_TOUCH\_PAD\textcolor{red}{\it{n}}\_REG (\textcolor{red}{\it{n}}: 0-7)、RTCIO\_TOUCH\_PAD\textcolor{red}{\it{m}}\_REG (\textcolor{red}{\it{m}} = 8, 9)
    \end{itemize}

    \item 章节 \ref{mod:spi} \textit{\nameref{mod:spi}}：标记寄存器 \hyperlink{regdesc:SPICTRL2REG}{SPI\_CTRL2\_REG} 中的 bit12 \verb+~+ bit15 为保留位 %doc-7416

    \item 章节 \ref{mod:sdioslave} \textit{\nameref{mod:sdioslave}}：修正几处笔误

    \item 章节 \ref{mod:ethernetmac} \textit{\nameref{mod:ethernetmac}}：删除寄存器 EMACCSTATUS\_REG，将其标记为 reserved %doc-7447

    \item 章节 \ref{mod:i2s} \textit{\nameref{mod:i2s}}：更新 \ref{DMA Interrupts} 小节中的中断描述 %DOC-6346

    \item 章节 \ref{mod:ledpwm} \textit{\nameref{mod:ledpwm}}：更新 \ref{LEDCchannel} 中关于不支持暂停 LEDC 渐变的说明，并更新表 \ref{tab:ledc-common-freq-res} 中的最低精度  %DOC-7304 %DIG-133

    \item 章节 \ref{mod:timg} \textit{\nameref{mod:timg}}：更新 TIMG\textcolor{red}{\it{n}}\_WDT\_CLK\_PRESCALE 的描述 %DOC-5697

    \item 章节 \ref{mod:sensor} \textit{\nameref{mod:sensor}}：更新寄存器 \hyperlink{regdesc:APBSARADCFSMREG}{APB\_SARADC\_FSM\_REG} %doc-7537

    \item 章节 \ref{mod:lowpowm} \textit{\nameref{mod:lowpowm}}：更新寄存器 \hyperref[regdesc:RTCCNTLWDTWPROTECTREG]{RTC\_CNTL\_WDTWPROTECT\_REG} 的描述 %DOC-6760

\end{itemize}
\vspace{-1em}
\\\hline
2023.07 & v5.0 &
\vspace{0.5em}
\begin{itemize}
    \item 全文更新 APB\_CTRL 寄存器前缀为 SYSCON   %DOC-5139
    \item 章节 \ref{cache} \textit{\nameref{cache}}：说明 cache 的块大小为 32 字节 %doc-4952
    \item 章节 \ref{mod:iomuxgpio} \textit{\nameref{mod:iomuxgpio}}：将原来的 RTCIO\_TOUCH\_PAD\textcolor{red}{\it{n}}\_REG (\textcolor{red}{\it{n}}: 0-9) 寄存器分成两个：\hyperref[regdesc:RTCIOTOUCHPADnREG]{RTCIO\_TOUCH\_PAD\textcolor{red}{\it{n}}\_REG (\textcolor{red}{\it{n}}: 0-7)} 和 \hyperref[regdesc:RTCIOTOUCHPADmREG]{RTCIO\_TOUCH\_PAD\textcolor{red}{\it{m}}\_REG (\textcolor{red}{\it{m}}: 8-9)}；在 \hyperref[regdesc:RTCIOTOUCHPADnREG]{RTCIO\_TOUCH\_PAD\textcolor{red}{\it{n}}\_REG (\textcolor{red}{\it{n}}: 0-7)} 中增加 bit [27-31]、[12-16] %doc-4910
    \item 章节 \ref{sec:rtc_gpio_analog_func} \textit{\nameref{sec:rtc_gpio_analog_func}}：修改描述
    \item 章节 \ref{mod:dportreg} \textit{\nameref{mod:dportreg}}：增加两个寄存器 \hyperref[regdesc:PROCACHECTRL1REG]{DPORT\_PRO\_CACHE\_CTRL1\_REG}、\hyperref[regdesc:APPCACHECTRL1REG]{DPORT\_APP\_CACHE\_CTRL1\_REG} 的描述  %doc-4728
    \item 章节 \ref{sec:peri-clock-gating-reset} \textit{\nameref{sec:peri-clock-gating-reset}}：增加说明“使用外设前必须先把时钟打开、复位释放” %DOC-4939
    \item 章节 \ref{sec:dma-spi-dma} \textit{\nameref{sec:dma-spi-dma}}：将 SPI DMA 一次接收/发送的数据长度改为“应当与 4 字节对齐” %DIG-154
    \item 章节 \ref{mod:ethernetmac} \textit{\nameref{mod:ethernetmac}}：删除时间戳/PTP 功能相关内容，原因是硬件暂不支持此功能 %DOC-4250
    \item 章节 \ref{mod:pcnt} \textit{\nameref{mod:pcnt}}：增加 \hyperref[regdesc:PCNTCTRLREG]{PCNT\_CLK\_EN} 位的描述 % DOC-5671
\end{itemize}
\vspace{-1em}
\\\hline
2023.04 & v4.9 &
\vspace{0.5em}
        \begin{itemize}
            \item 根据 \href{https://www.espressif.com/sites/default/files/pcn_downloads/PCN20221202%20Remove%20Hall%20Sensor%20from%20ESP32%20Series%20of%20Documentation.pdf}{PCN20221202} 删除霍尔传感器相关内容、寄存器、信号等
            \item 将全文出现的 PLL\_D2\_CLK 改为 PLL\_F160M\_CLK
            \item 章节 \ref{mod:iomuxgpio} \textit{\nameref{mod:iomuxgpio}}：在表 \ref{table:peripheral_signals} 中增加 TWAI 信号

            \item 在章节 \ref{mod:uart} \textit{\nameref{mod:uart}} 新增终止状态 (break condition) 的相关描述，更新停止位最大位数和相关描述%DOC-3633
            \item 在章节 \ref{mod:ledpwm} \textit{\nameref{mod:ledpwm}} 新增占空比精度计算公式，并更新表 \textit{\nameref{tab:ledc-common-freq-res}} %DOC-3480
            \item 章节 \ref{mod:sensor} \textit{\nameref{mod:sensor}}：
            \begin{itemize}
                \item 在小节 \ref{sec:on_chip_sensor_features} \textit{\nameref{sec:on_chip_sensor_features}} 中增加应用限制的说明
                \item 删除内部信号 vdd33、pa\_pkdet1、pa\_pkdet2
            \end{itemize}
            \item 在章节 \ref{mod:lowpowm} \textit{\nameref{mod:lowpowm}} 中增加有关“拒绝睡眠”的描述 %DOC-4606
        \end{itemize}

        \vspace{-1em}
    \\\hline
        2022.12   & v4.8 &
        \vspace{0.5em}
        更新以下章节：
        \begin{itemize}
            \item 在章节 \ref{mod:resclk} \textit{\nameref{mod:resclk}} 中更新小节 \ref{sec:ref_tick}，新增寄存器小节 \ref{sec:resclk-reg-summ} 和 \ref{sec:resclk-reg-desc}
            \item 更新章节 \ref{mod:ethernetmac} \textit{\nameref{mod:ethernetmac}} 中 \hyperref[fielddesc:MIICSRCLK]{MIICSRCLK} 和 \hyperref[fielddesc:MIIBUSY]{MIIBUSY} 字段的描述
            \item 在章节 \ref{mod:uart} \textit{\nameref{mod:uart}} 新增 \hyperref[regdesc:UARTMEMTXSTATUSREG]{UART\_MEM\_TX\_STATUS\_REG} 和 \hyperref[regdesc:UARTMEMRXSTATUSREG]{UART\_MEM\_RX\_STATUS\_REG} 寄存器的描述 %DOC-1718
            \item 更新章节 \ref{mod:rng} \textit{\nameref{mod:rng}} 中的时钟名称
            \item 更新章节 \ref{mod:sensor} \textit{\nameref{mod:sensor}} 中的表 \ref{table:adc inputs}
            \item 更新章节 \ref{mod:lowpowm} \textit{\nameref{mod:lowpowm}}  中的时钟名称，并增加有关 \hyperref[fielddesc:RTCCNTLMAINSTATEINIDLE]{RTC\_CNTL\_MAIN\_STATE\_IN\_IDLE} 的描述
            \item 将章节 \ref{mod:efuse} \textit{\nameref{mod:efuse}} 中的 RTC8M\_CLK 重命名为 RC\_FAST\_CLK
            \item 在所有章节的寄存器小节增加了有关“相对地址”的说明
            \vspace{-1em}
        \end{itemize}
        \\\hline
        2022.08   & v4.7 &
        \vspace{0.5em}
    \begin{itemize}
        \item 章节 \ref{mod:sysmem} \textit{\nameref{mod:sysmem}} ：更新表 \ref{table:Peripheral_Address_Mapping} %doc-2386
        \item 章节 \ref{mod:iomuxgpio} \textit{\nameref{mod:iomuxgpio}}：
    \begin{itemize}
        \item 修正寄存器列表中的一处笔误 %doc-2874
        \item 更新第 \ref{sec:inputconfig} 小节中的描述 %doc-2880
        \item 更新对 GPIO\_PIN\regindex{n}\_INT\_ENA 和 IO\_MUX\_PIN\_CTRL 的描述%doc-3018 %doc-2081
    \end{itemize}
    \item 章节 \ref{mod:dma} \textit{\nameref{mod:dma}} ：在 \ref{sec:UDMA} \textit{\nameref{sec:UDMA}} 小节新增一条关于接收链表描述符的说明 %dig-82
    \item 章节 \ref{mod:spi} \textit{\nameref{mod:spi}} ：更新 SPI\_TRANS\_DONE、SPI\_SLV\_WR\_STA\_DONE、SPI\_SLV\_RD\_STA\_DONE、SPI\_SLV\_WR\_BUF\_DONE、SPI\_SLV\_RD\_BUF\_DONE 的描述 %doc-3016
    \item 更新时钟名称：
    \begin{itemize}
        \item RTC8M\_CLK 更新为 RC\_FAST\_CLK
        \item RTC8M\_D256\_CLK 更新为 RC\_FAST\_DIV\_CLK
        \item RTC\_CLK 更新为 RC\_SLOW\_CLK
        \item SLOW\_CLK 更新为 RTC\_SLOW\_CLK
        \item FAST\_CLK 更新为 RTC\_FAST\_CLK
    \end{itemize}%doc-3141
    \item 章节 \ref{mod:i2s} \textit{\nameref{mod:i2s}} ：修正寄存器列表中的一处笔误 %doc-2919
    \item 章节 \ref{mod:timg} \textit{\nameref{mod:timg}} ：更新寄存器列表 %doc-2822
    \item 章节 \ref{mod:sensor} \textit{\nameref{mod:sensor}} ：新增两条说明 %doc-1694
        \end{itemize}
    \vspace{-1em}
    \\\hline
    2021.11 & v4.6 &
    更新章节 \ref{mod:sysmem} \textit{\nameref{mod:sysmem}} 中的表 \ref{table:Peripheral_Address_Mapping} %doc-1520/doc-1930

    更新章节 \ref{mod:intmtrx} \textit{\nameref{mod:intmtrx}} 中的表 \ref{table:PeripheralInterruptConfigStatusSource}  %doc-2011

    更新章节 \ref{mod:resclk} \textit{\nameref{mod:resclk}} 中的图 \ref{figure:Reset}  %doc-2091

    更新章节 \ref{mod:iomuxgpio} \textit{\nameref{mod:iomuxgpio}} 中的表 \ref{table:peripheral_signals} 和对 \hyperref[regdesc:IOMUXPINCTRL]{IO\_MUX\_PIN\_CTRL} 的描述  %doc-2311, %doc-2081

    全面更新章节 \ref{mod:dportreg} \textit{\nameref{mod:dportreg}}%doc-1520

    在章节 \ref{mod:i2s} \textit{\nameref{mod:i2s}} 的表 \ref{table:I2S 信号总线描述} 下方新增一条说明 %doc-1846

    更新章节 \ref{mod:sdioslave} \textit{\nameref{mod:sdioslave}} 中的 \ref{SDIO slave timing} 小节和寄存器 \hyperref[regdesc:SLCHOSTCONFREG]{SLCHOST\_CONF\_REG} %doc-1985, %doc-1740

    更新章节 \ref{mod:extmemencr} \textit{\nameref{mod:extmemencr}} 的描述 %doc-2305

    更新章节 \ref{mod:sensor} \textit{\nameref{mod:sensor}} 中的 \ref{sec:structure} 小节 %doc-2017

    在章节 \ref{mod:iomuxgpio}、章节 \ref{mod:spi}、章节 \ref{mod:i2s}、章节 \ref{mod:rmt}、章节 \ref{mod:sensor} 和章节 \ref{mod:ulp} 中增加对寄存器地址的说明 %DGAP-42, doc-1120

    翻译多个章节中的寄存器列表 %doc-478

    更新词汇列表 %doc-2115

    修正笔误 %doc-1589 %doc-2146
    \\\hline
2021.07 & v4.5 &
增加章节 \ref{section:PHY_Timing} %DOC-IDFGH-1673

在表 \ref{table:rtcmux pads} 下方增加一条说明 %DOC-1644

在章节 \ref{mod:ledpwm} \textit{\nameref{mod:ledpwm}} 更新表 \ref{tab:ledc-common-freq-res} %DOC-1420

更新章节 \ref{audiopll} 的描述 %DOC-1494

将章节 \ref{section:Ethernet Mac Register Summary} 中的 Ethernet MAC 基地址更新为 0x3FF6\_9000 %DOC-1514

在章节 \ref{mod:uart} \textit{\nameref{mod:uart}} 更新 UART\_SW\_RTS、UART\_RX\_FLOW\_THRHD 字段的描述 % DOC-1499, DOC-1515

在章节 \ref{mod:twai} \textit{\nameref{mod:twai}} 更新特性描述 %DOC-1646

更新章节 \ref{mod:rng} \textit{\nameref{mod:rng}} % DOC-1257

更新章节 \ref{st instruction} 的描述 %DOC-1355

在章节 \ref{mod:dportreg} \textit{\nameref{mod:dportreg}} 和章节 \ref{mod:intmtrx} \textit{\nameref{mod:intmtrx}} 中添加中断矩阵寄存器的基地址信息 % DOC-1382

更新章节 \ref{mod:iomuxgpio} \textit{\nameref{mod:iomuxgpio}} 中管脚功能编号，编号从 Function0 开始  %DOC-1523

删除章节 \ref{sec:baudrate-detect} 中重复的内容 % DOC-1398

修正表 \ref{table:PeripheralInterruptConfigStatusSource} 和章节 \ref{sec:UDMA} 中的笔误 %DOC-1328, DOC-1621
\\\hline

2021.03 & v4.4 &
更新章节 \ref{mod:mcpwm}：\textit{\nameref{mod:mcpwm}} 中寄存器 \hyperref[regdesc:PWMTIMER0SYNCREG]{PWM\_TIMER0\_SYNC\_REG} \verb+~+ \hyperref[regdesc:PWMTIMER2SYNCREG]{PWM\_TIMER2\_SYNC\_REG} 的描述

更新章节 \ref{mod:timg}：\textit{\nameref{mod:timg}} 中 \nameref{regdesc:TIMGRTCCALICFGREG} 寄存器和 \nameref{regdesc:TIMGRTCCALICFG1REG} 寄存器的描述

更新章节 \ref{mod:lowpowm}：\textit{\nameref{mod:lowpowm}} 中 \nameref{regdesc:RTCCNTLWDTCONFIG0REG} 寄存器的描述

更新章节 \ref{mod:ethernetmac}：\textit{\nameref{mod:ethernetmac}} 中对 \nameref{fielddesc:MIICSRCLK} 的描述

更新章节 \ref{mod:iomuxgpio}：\textit{\nameref{mod:iomuxgpio}} 中对 \nameref{regdesc:FUNWPU} 和 \nameref{regdesc:FUNWPD} 的描述

更新章节 \ref{mod:twai}：\textit{\nameref{mod:twai}} 中所用的商标符号

更新章节 \ref{mod:lowpowm}：\textit{\nameref{mod:lowpowm}} 中 \hyperlink{regdesc:RTCCNTLULPCPTIMERREG}{RTC\_CNTL\_ULP\_CP\_TIMER\_REG} 寄存器的描述

将文档写作规范章节重命名为词汇列表，并移至文档结尾

更新表 \ref{table:iomux_pads}：\textit{\nameref{table:iomux_pads}}下方的说明
\\\hline
2020.09 & v4.3 &
新增章节 \ref{mod:twai} \textit{\nameref{mod:twai}}

在章节 \hyperref[mod:efuse]{\it{eFuse 控制器}} 中新增 uart\_download\_dis 系统参数的相关信息。

在章节 \ref{mod:rng} \textit{\nameref{mod:rng}} 中新增 \ref{sec:rng-programming} \textit{\nameref{sec:rng-programming}} 小节，并更新部分描述。

更新对寄存器 \hyperref[regdesc:SPIADDRREG]{SPI\_ADDR\_REG} 和 \hyperref[regdesc:SPISLVWRSTATUSREG]{SPI\_SLV\_WR\_STATUS\_REG} 的描述
\\\hline
2020.06 & v4.2 &
新增文档写作规范章节

更新章节\hyperref[mod:sysmem]{\it{系统和存储器}}：
\begin{itemize}
\item 在表 \ref{table:Peripheral_Address_Mapping} 下方添加关于 DPORT 和 AHB 地址空间的说明
\end{itemize}

更新章节\hyperref[mod:resclk]{\it{复位和时钟}}：
\begin{itemize}
\item 更新表 \ref{table:cpu clk}：增加 PLL\_CLK 频率为 480 MHz，CPU\_CLK 频率为 240 MHz 的描述
\item 更新表 \ref{table:apb_clk_src}：将 CPU\_CLK 源为 PLL\_CLK 时 APB\_CLK 的频率改为 80 MHz
\end{itemize}

更新章节 \hyperref[mod:iomuxgpio]{\it{IO\_MUX 和 GPIO 交换矩阵}}：
\begin{itemize}
\item 修正章节 \ref{direct io via iomux} 中的一处笔误：对于输入信号，必须清零 SIG\_IN\_SEL 寄存器，直接将输入信号输出到外设。
\item 修改表 \ref{table:iomux_pads} 中 MTCK、MTMS、GPIO27 的 reset 状态
\item 更新寄存器 \hyperref[FUNDRV]{FUN\_DRV} 的描述
\end{itemize}

更新章节 \hyperref[mod:i2s]{\it{I2S}}:
\begin{itemize}
\item 更新章节 \ref{section:Philips Standard}
\end{itemize}

更新章节 \hyperref[mod:uart]{UART 控制器}：
\begin{itemize}
\item 更新章节 \ref{section:UART RAM}
\item 更新寄存器 \hyperref[regdesc:UARTFIFOREG]{UART\_FIFO\_REG} 和 \hyperref[UARTRXTOUTTHRHD]{UART\_RX\_TOUT\_THRHD} 的描述
\end{itemize}

更新章节 \hyperref[mod:ledpwm]{\it{LED\_PWM}}：
\begin{itemize}
\item 增加表 \ref{tab:ledc-common-freq-res}
\end{itemize}

更新章节\hyperref[mod:mcpwm]{\it{电机控制脉宽调制器 (MCPWM)}}：
\begin{itemize}
\item 修正在递增计数模式、递减计数模式、递增-递减循环模式下 PWM 的周期值
\end{itemize}

更新章节 \hyperref[mod:pcnt]{\it{PULSE\_CNT}}：
\begin{itemize}
\item 增加寄存器 \hyperref[regdesc:PCNTUnSTATUSREG]{PCNT\_U\textcolor{red}{\it{n}}\_STATUS\_REG} 的描述
\end{itemize}

更新章节 \hyperref[mod:efuse]{\it{eFuse 控制器}}：
\begin{itemize}
\item 将两个系统参数“32pad”和“chip\_version”合并成一个：pkg\_version
\item 更新寄存器 \hyperref[regdesc:EFUSERDCHIPVERPKG]{EFUSE\_RD\_CHIP\_VER\_PKG} 和 \hyperref[regdesc:EFUSECHIPVERPKG]{EFUSE\_CHIP\_VER\_PKG} 的描述
\end{itemize}

更新章节\hyperref[mod:sensor]{\it{片上传感器与模拟信号处理}}：
\begin{itemize}
\item 修正几处笔误
\end{itemize}

更新章节\hyperref[mod:ulp]{\it{超低功耗协处理器}}：
\begin{itemize}
\item 更新章节 \ref{subsection:REG_RD} 和 \ref{subsection:REG_WR} 中的描述
\item 修正几处笔误
\end{itemize}

更新章节\hyperref[mod:lowpowm]{低功耗管理}：
\begin{itemize}
\item 增加寄存器 \hyperref[regdesc:RTCCNTLWDTCONFIG0REG]{RTC\_CNTL\_WDTCONFIG0\_REG} 的描述
\item 修改寄存器 \hyperref[regdesc:RTCCNTLWDTCONFIGnREG]{RTC\_CNTL\_WDTCONFIG\textcolor{red}{\it{n}}\_REG} 的描述
\end{itemize}
\vspace{-1em}
\\\hline

2019.11 &v4.1 &
更新章节 \hyperref[mod:iomuxgpio]{\it{IO\_MUX 和 GPIO 交换矩阵}}：
\begin{itemize}
\item 更新表 \ref{table:rtcmux pads}；
\item 在寄存器 \hyperref[regdesc:RTCIOTOUCHPADnREG]{{RTCIO\_TOUCH\_PAD\textcolor{red}{\it{n}}\_REG}} 描述中增加 RTCIO\_TOUCH\_PAD\textcolor{red}{\it{n}}\_FUN\_SEL 及其说明；
\end{itemize}

更新章节 \hyperref[mod:spi]{\it{SPI}}：
\begin{itemize}
\item 修正表 \ref{Section:SPIREGList} 中 SPI2、SPI3 的错误地址；
\end{itemize}

更新章节 \hyperref[mod:i2c]{\it{I2C 控制器}}：
\begin{itemize}
\item 删除 I2C\_SLAVE\_TRAN\_COMP\_INT 中断；
\end{itemize}

更新章节 \hyperref[mod:i2s]{\it{I2S}}：
\begin{itemize}
\item 在表 \ref{table:I2S 信号总线描述} 下增加一条说明；
\end{itemize}

更新章节 \hyperref[mod:uart]{\it{UART 控制器}}：
\begin{itemize}
\item 修正 \hyperref[regdesc:UARTFORCEXOFF]{UART\_FORCE\_XOFF} 和 \hyperref[regdesc:UARTFORCEXON]{UART\_FORCE\_XON} 的寄存器描述；
\item 修正 \hyperref[regdesc:UARTSWFCCONFREG]{ART\_SWFC\_CONF\_REG} 的寄存器描述；
\end{itemize}

更新章节\hyperref[mod:rmt]{\it{红外遥控}}：
\begin{itemize}
\item 更新图 \ref{Figure:RMT2-1}；
\end{itemize}

更新章节 \hyperref[mod:pcnt]{\it{PULSE\_CNT}}：
\begin{itemize}
\item 更新图 \ref{Figure:pulse_cnt_architecture}；
\item 修正寄存器 \hyperref[regdesc:PCNTUnCONF0REG]{PCNT\_U\textit{\textcolor{red}{\it{n}}}\_CONF0\_REG} 描述里的笔误；
\end{itemize}

更新章节 \hyperref[mod:efuse]{\it{eFuse 控制器}}：
\begin{itemize}
\item 新增并修改 8 个系统参数；更新相应寄存器；
\item 更新表 \ref{table:Timing_Configure} 中的寄存器配置值；
\item 将系统参数 flash\_crypt\_cnt 的位宽改为 7 位；
\end{itemize}

更新章节 \hyperref[mod:mmumpu]{\it{PID/MPU/MMU}}：
\begin{itemize}
\item 在表格 \ref{table:MPU_for_DMA} 下方增加一条说明；
\end{itemize}

更新章节\hyperref[mod:sensor]{\it{片上传感器与模拟信号处理}}：
\begin{itemize}
\item 修正 \hyperref[SENSSAR2BITWIDTH]{SENS\_SAR2\_BIT\_WIDTH} 和 \hyperref[SENSSAR1BITWIDTH]{SENS\_SAR1\_BIT\_WIDTH} 的寄存器描述；
\end{itemize}

更新章节\hyperref[mod:ulp]{\it{超低功耗协处理器}}：
\begin{itemize}
\item 删除 TSENS 指令；
\item 修正 \hyperref[subsection:REG_WR]{REG\_WR} 的 OpCode；
\item 更新章节 \ref{Connecting I2C Signals}；
\item 修正 \hyperref[RTCI2CRXLSBFIRST]{RTC\_I2C\_RX\_LSB\_FIRST} 和 \hyperref[RTCI2CTXLSBFIRST]{RTC\_I2C\_TX\_LSB\_FIRST} 的寄存器描述；
\item 删除 RTC\_I2C\_SLAVE\_TRAN\_COMP\_INT\_ENA 和 RTC\_I2C\_SLAVE\_TRAN\_COMP\_INT\_ST 的寄存器描述；
\end{itemize}

更新章节\hyperref[mod:lowpowm]{\it{低功耗管理}}：
\begin{itemize}
\item 修正寄存器 \hyperref[RTCCNTLDBROWNOUTTHRES]{RTC\_CNTL\_DBROWN\_OUT\_THRES} 的默认值和描述；
\item 修正寄存器 \hyperref[RTCCNTLBROWNOUTCLOSEFLASHENA]{RTC\_CNTL\_BROWN\_OUT\_CLOSE\_FLASH\_ENA} 的描述；
\end{itemize}

增加\href{https://www.espressif.com/zh-hans/company/contact/documentation_feedback?docId=3641&sections=}{文档反馈链接}。
\\\hline

2018.12&v4.0&更新章节 \hyperref[mod:iomuxgpio]{IO\_MUX 和 GPIO 交换矩阵} 中几处寄存器名称，以跟头文件保持一致；

更新章节 \ref{mod:spi} SPI：
\begin{itemize}
    \item 更新章节 \ref{GP-SPI}；
    \item 更新章节 \ref{Communication Format of Parallel QSPI}；
    \item 更新章节 \ref{Section:SPIREG}；
\end{itemize}

更新章节 \ref{mod:uart} UART 控制器：
\begin{itemize}
    \item 将 UART 支持的停止位修正为 1/1.5/2/3 个；
    \item 在章节 \ref{UARTArchitecture} 末段增加一则说明；
    \item 更新寄存器 \hyperref[UARTDL0EN]{UART\_DL0\_EN} 的描述；
\end{itemize}
在\hyperref[mod:lowpowm]{低功耗管理}一章中\hyperref[WakeupSource]{唤醒源}的末段增加一则说明。\\\hline
2018.10&v3.9& 更新章节 \ref{mod:i2c}：I2C 控制器中图 \ref{Figure:I2C_bus_timing}：I2C 时序图。\\\hline
2018.09&v3.8&更新寄存器 \hyperref[TIMGnTxALARMEN]{TIMG\textcolor{red}{\it{n}}\_T\textcolor{red}{\it{x}}\_ALARM\_EN} 的描述；

在章节 \ref{ULP Program Execution} 中增加关于 ULP 协处理器唤醒时间的描述。\\\hline
2018.08&v3.7&更新寄存器 \hyperref[UARTRXGAPTOUT]{UART\_RX\_GAP\_TOUT} 的描述。\\\hline
\multirow{2}*{2018.08}&\multirow{2}*{v3.6}&更新章节 \ref{jumpstorelativeaddress} 中跳转条件；\\
&&更新寄存器 \hyperref[UARTACTIVETHRESHOLD]{UART\_ACTIVE\_THRESHOLD} 的描述。
\\\hline
2018.07&v3.5&
更新章节 \ref{mod:rmt} RMT：
\begin{itemize}
\item 更新章节 \ref{RMTRAM}：RMT RAM 中 RAM 的起始地址；
\item 修正几处 RMT 寄存器地址错误；
\item 更新寄存器 \hyperref[regdesc:RMTAPBCONFREG]{RMT\_APB\_CONF\_REG} 的描述。
\end{itemize}
更新寄存器 \hyperref[UARTRXTOUTTHRHD]{UART\_RX\_TOUT\_THRHD}、\hyperref[UARTRXFIFOFULLINTCLR]{UART\_RXFIFO\_FULL\_INT\_CLR}、\hyperref[UARTRXFIFOFULLINTCLR]{UART\_RXFIFO\_FULL\_INT\_CLR} 的描述。
\\\hline
\multirow{5}*{2018.06}&\multirow{5}*{v3.4}&
更新章节 \ref{I/O Pad Power Supplies} ESP32 I/O Pad 供电中的图；\\
&&更新章节 \ref{I2C Bus Timing} I2C 总线时序；\\
&&在章节 \ref{LEDCchannel} LEDC 通道中增加说明；\\
&& 更新章节 \ref{Watchpoints} PULSE\_CNT 观察点中的“最大计数值”；\\
&&删除有关温度传感器和超低噪声前置模拟放大器相关的内容。\\\hline
2018.05 & v3.3 & 更新 \hyperref[mod:lowpowm]{低功耗管理} 章中 \hyperref[register summary-30]{寄存器列表} 节和 \hyperref[registers-30]{寄存器} 节中的寄存器地址。 \\\hline
\multirow{9}*{2018.04}&\multirow{9}*{v3.2}& 更新图 \ref{figure:CMD53内容} CMD53 内容。\\
&& 在章节 \hyperref[mod:ethernetmac]{Ethernet MAC} 中增加以下 7 个寄存器：  \begin{itemize} \item\hyperref[regdesc:DMAOPERATION_MODEREG]{DMAOPERATION\_MODE\_REG};
\item\hyperref[regdesc:DMAINENREG]{DMAIN\_EN\_REG},
\item\hyperref[regdesc:DMAMISSEDFRREG]{DMAMISSEDFR\_REG}, \item\hyperref[regdesc:PMT_RWUFFRREG]{PMT\_RWUFFR\_REG},
\item\hyperref[regdesc:PMT_CSRREG]{PMT\_CSR\_REG},
\item\hyperref[regdesc:EMACLPI_CSRREG]{EMACLPI\_CSR\_REG}, and
\item\hyperref[regdesc:EMACLPITIMERSCONTROLREG]{EMACLPITIMERSCONTROL\_REG}.\vspace{-1em}\end{itemize}\\\hline

2018.04 & v3.1 & 更新图 \ref{Figure:RMT2-1} RMT 架构；\\
&&在章节 \ref{PadHold} 中增加说明；\\
&&在章节 \ref{regdesc:RTCIODIGPADHOLDREG} 中增加对寄存器位的描述；\\\hline
2018.03& v3.0 & 更新章节 \ref{st instruction} 中 \hyperref[figure:ST]{ST 指令的定义图}；\\ &&
 在章节 \ref{section:Ethernet Mac Register Summary} 和 \ref{ethernetreg} 中增加寄存器 \hyperref[regdesc:EMACADDR2HIGHREG]{EMACADDR2HIGH\_REG} 到 \hyperref[regdesc:EMACADDR7LOWREG]{EMACADDR7LOW\_REG} 说明。\\\hline
2018.02 & v2.9 & 更新章节 \ref{sec:inputconfig}、\ref{IOMUXSimpleGPIOInput}、\ref{IOMUXPeripheralOutput}；\\
&&在章节 \hyperref[I2SRegister]{I2S 寄存器} 中增加 I2S\_FIFO\_WR\_REG 和 I2S\_FIFO\_RD\_REG 寄存器。\\\hline
\multirow{2}*{2018.01}&\multirow{2}*{v2.8}&增加章节\hyperref[EMAC]{以太网 MAC}。\\ &&在章节 \hyperref[mod:efuse]{eFuse 控制器} 中增加系统参数 \hyperref[BLK3]{BLK3\_part\_reserve} 的描述。\\\hline
\multirow{3}*{2017.12}&\multirow{3}*{v2.7}&在章节 \hyperref[mod:sysmem]{系统和存储器} 中增加小节 \hyperref[cache]{Cache}；\\
&&更新章节 \hyperref[ledc_timers]{分频器} 及 \hyperref[mod:ledpwm]{LED\_PWM} 中多个寄存器名称；\\
&& 更新章节 \hyperref[mod:efuse]{eFuse 控制器} 中寄存器 \hyperref[console]{console\_debug\_disable} 的描述。\\\hline
2017.11&v2.6&更新章节\hyperref[mod:rmt]{红外遥控}：
\begin{itemize}
\item 更新图 \ref{Figure:RMT2-1} RMT 架构；
\item 更新章节 \hyperref[RMTRAM]{RMT RAM}；
\item 更新章节 \hyperref[rmt:transmitter]{发射器}；
\item 更新中断 \hyperref[int:RMTCHnTXTHREVENTINT]{RMT\_CH\textcolor{red}{\it{n}}\_TX\_THR\_EVENT\_INT} 的描述。
\end{itemize}
在章节 \hyperref[section:UART RAM]{UART RAM} 和寄存器 \hyperref[regdesc:UARTCONF0REG]{UART\_CONF0\_REG} 中增加说明。
\\\hline
\multirow{4}*{2017.11}&\multirow{4}*{v2.5}& 更新章节 \hyperref[Section:SPIREGList]{ SPI 寄存器列表}中寄存器 SPI\_CTRL\_REG 的地址描述；\\
&&在章节 SD/MMC 主机控制器中增加\hyperref[sdmmc_phase_sel]{时钟相位选择}，增加寄存器 \hyperref[regdesc:CLKEDGESEL]{CLK\_EDGE\_SEL} 的说明；\\
&&关于 \hyperref[mod:i2c]{I2C 控制器}章节的重大修订。\\\hline
\multirow{5}*{2017.09}&\multirow{5}*{v2.4}&在章节 \hyperref[mod:sdioslave]{SDIO 从机} 中增加寄存器 \hyperref[regdesc:SLC0HOSTTOKENRDATA]{SLC0HOST\_TOKEN\_RDATA} 的描述；\\
&&在章节 \hyperref[The Clock of I2S Module]{I2S 模块时钟}中增加注意事项；\\
&&在章节 GP-SPI 主机模式中增加说明；\\
&&增加章节 \hyperref[mod:dportreg]{DPort 寄存器}；\\
&&增加章节 \hyperref[mod:dma]{DMA 控制器}。\\\hline
2017.08&v2.3&增加章节 \hyperref[mod:extmemencr]{Flash 加密与解密}。\\\hline
2017.07&v2.2&增加章节 \hyperref[mod:lowpowm]{低功耗管理}。\\\hline
2017.07&v2.1&更新章节 \hyperref[section:4.12]{IO\_MUX 和 GPIO 交换矩阵}中 GPIO 配置/数据寄存器和 GPIO RTC 功能配置寄存器的地址；

增加章节 \hyperref[mod:pid]{PID 控制器}。\\\hline
2017.07&v2.0&增加章节 \hyperref[mod:sdioslave]{SDIO 从机}。\\\hline
\multirow{2}*{2017.06} &\multirow{2}*{v1.9}&更新章节 \hyperref[mod:iomuxgpio]{IO\_MUX 和 GPIO 交换矩阵}；\\
&& 增加章节 \hyperref[mod:mcpwm]{电机控制脉宽调制器（MCPWM）}。\\\hline
\multirow{3}*{2017.06} & \multirow{3}*{v1.8} & 在章节 \hyperref[mod:i2s]{I2S} 中增加寄存器 \hyperref[regdesc:I2SSTATEREG]{I2S\_STATE\_REG} 的描述；\\
&&更新章节 \hyperref[mod:iomuxgpio]{IO\_MUX 和 GPIO 交换矩阵}；\\
&&增加章节 \hyperref[mod:ulp]{超低功耗协处理器}。
\\\hline
\multirow{5}*{2017.05} &\multirow{5}*{v1.7} & 增加章节 \hyperref[mod:sensor]{片上传感器与模拟信号处理}；\\
&&增加章节 \hyperref[audiopll]{音频 PLL}；\\
&&更新章节 \hyperref[15.4]{eFuse 控制器寄存器列表}；\\
&&更新章节 \hyperref[I2SPDM]{I2S PDM 模式}和 \hyperref[LCDMODE]{LCD 模式}；\\
&& 更新章节：GP-SPI 从机支持的通信格式。
\\\hline
\multirow{3}*{2017.03} &\multirow{3}*{v1.6} & 增加章节 \hyperref[mod:sdhost]{SD/MMC 主机控制器}；\\
&&在章节 \hyperref[mod:iomuxgpio]{IO\_MUX 和 GPIO 交换矩阵}中增加寄存器 \hyperref[regdesc:IOMUXPINCTRL]{IO\_MUX\_PIN\_CTRL} 的描述。
\\\hline
2017.03 & v1.5 & 增加章节 \hyperref[mod:i2s]{I2S}。
\\\hline
\multirow{2}*{2017.01} & \multirow{2}*{v1.4} & 增加章节 \hyperref[mod:spi]{SPI}；\\
&&增加章节 \hyperref[mod:uart]{UART 控制器}。\\\hline
\multirow{4}*{2016.12} & \multirow{4}*{v1.3} & 增加章节 \hyperref[mod:efuse]{eFuse 控制器}；\\
&&增加章节 \hyperref[mod:rsa]{RSA 加速器}；\\
&& 增加章节 \hyperref[mod:rng]{随机数发生器}；\\
&& 更新章节 \hyperref[5.3.7]{I2C 控制器中断} 和 \hyperref[5.5]{I2C 控制器寄存器}。
\\\hline
\multirow{3}*{2016.11} & \multirow{3}*{v1.2} & 增加章节 \hyperref[mod:mmumpu]{PID/MPU/MMU}；\\
&&更新章节 \hyperref[section:4.12]{IO\_MUX 和 GPIO 交换矩阵寄存器列表}；\\
&& 更新章节 \hyperref[section:6.3]{LED\_PWM 寄存器列表}。
\\\hline
2016.09 & v1.1 & 增加章节 \hyperref[mod:i2c]{I2C 控制器}。
\\\hline
2016.08 & v1.0 & 首次发布。
\\\hline

\end{longtable}
